\section*{Bab \ref{ch:g11:fossil-fuel-geology} --- Geologi Batu Bara, Minyak Bumi, dan Gas}

\begin{solution}{exo:g11:fossil-fuel-geology:1}
Plankton mati mengendap di dasar laut yang miskin oksigen; bahan organik
menjadi \omterm{def:g11:fossil-fuel-geology:source-rock}{kerogen} di dalam lumpur; penimbunan di bawah sedimen yang lebih muda;
\omterm{def:g11:fossil-fuel-geology:oil-window}{pematangan} dalam \omterm{def:g11:fossil-fuel-geology:oil-window}{jendela minyak}; \omterm{def:g11:fossil-fuel-geology:migration}{migrasi}; minyak bumi berkumpul di bawah \omterm{def:g11:fossil-fuel-geology:reservoir}{batuan
penutup}.
\end{solution}

\begin{solution}{exo:g11:fossil-fuel-geology:2}
(1) \omterm{def:g11:fossil-fuel-geology:source-rock}{Batuan sumber}. (2) \omterm{def:g11:fossil-fuel-geology:reservoir}{Batuan reservoir}. (3) \omterm{def:g11:fossil-fuel-geology:reservoir}{Batuan penutup}. (4) \omterm{def:g11:fossil-fuel-geology:reservoir}{Batuan
penutup}. (5) \omterm{def:g11:fossil-fuel-geology:reservoir}{Batuan reservoir}.
\end{solution}

\begin{solution}{exo:g11:fossil-fuel-geology:3}
Puncak: $(60 - 10)/25 = \qty{2.0}{km}$; dasar: $(120 - 10)/25 = \qty{4.4}{km}$.
Tebal jendela itu $60/25 = \qty{2.4}{km}$.
\end{solution}

\begin{solution}{exo:g11:fossil-fuel-geology:4}
(1) Gambut, lignit, \omterm{def:g11:fossil-fuel-geology:coal}{batu bara} bituminus, antrasit. (2) Antrasit, 86 sampai
\qty{97}{\percent} karbon. (3) Antrasit, di luar sekitar \qty{5}{km} dan
\qty{150}{\degreeCelsius}.
\end{solution}

\begin{solution}{exo:g11:fossil-fuel-geology:5}
(1) Fluida-fluida memilah diri menurut massa jenisnya: gas paling ringan, lalu
minyak bumi, lalu air. (2) Minyak bumi dan gas naik karena lebih ringan daripada
air; hanya batuan yang tidak dapat dilaluinya yang menghentikannya. Batu lumpur
(atau serpih) dan garam.
\end{solution}

\begin{solution}{exo:g11:fossil-fuel-geology:6}
(1) A: $(60 - 15)/20 = \qty{2.25}{km}$ dan $(120 - 15)/20 = \qty{5.25}{km}$.
B: $45/40 \approx \qty{1.13}{km}$ dan $105/40 \approx \qty{2.63}{km}$.
(2) Pada \qty{2}{km}: di A, $15 + 40 = \qty{55}{\degreeCelsius}$, belum
matang, tanpa minyak bumi; di B, $15 + 80 = \qty{95}{\degreeCelsius}$, di dalam
\omterm{def:g11:fossil-fuel-geology:oil-window}{jendela minyak}, batuan itu menghasilkan minyak bumi. (3) Di cekungan A: minyak
bumi terbentuk di sana hanya di bawah \qty{2.25}{km}.
\end{solution}

\begin{solution}{exo:g11:fossil-fuel-geology:7}
(1) $11 \times 3 = \qty{33}{m}$ gambut. (2) $\qty{33000}{mm}$ dengan laju
\qty{1}{mm} setahun: sekitar \num{33000} tahun. (3) Gambut hanya bertahan di
tempat ia berada di bawah air yang tergenang, jauh dari jangkauan oksigen;
selagi gambut menumpuk, tanahnya harus turun dengan laju yang sama agar
permukaannya tetap terendam air.
\end{solution}

\begin{solution}{exo:g11:fossil-fuel-geology:8}
(1) $20 + 30 \times 1.2 = \qty{56}{\degreeCelsius}$;
$20 + 30 \times 2.3 = \qty{89}{\degreeCelsius}$;
$20 + 30 \times 4.0 = \qty{140}{\degreeCelsius}$.
(2) 120 juta tahun yang lalu, belum matang; 95 juta tahun yang lalu, di dalam
\omterm{def:g11:fossil-fuel-geology:oil-window}{jendela minyak}; 20 juta tahun yang lalu, di dalam \omterm{def:g11:fossil-fuel-geology:oil-window}{jendela gas}.
\end{solution}

\begin{solution}{exo:g11:fossil-fuel-geology:9}
(1) $2580 - 2400 = \qty{180}{m}$. (2) Gas mengisi
$2450 - 2400 = \qty{50}{m}$, minyak bumi $2580 - 2450 = \qty{130}{m}$.
(3) Keduanya melimpah keluar di titik limpah, di bawah kontur tertutup yang
terdalam, lalu bermigrasi lebih jauh ke atas, ke \omterm{def:g11:fossil-fuel-geology:trap}{perangkap} lain atau ke
permukaan.
\end{solution}

\begin{solution}{exo:g11:fossil-fuel-geology:10}
(1) Puncak jendela pada \qty{3}{km}: $(60 - 20)/3 \approx
\qty{13}{\degreeCelsius/km}$; dasarnya: $(120 - 20)/3 \approx
\qty{33}{\degreeCelsius/km}$. (2) Pada \qty{3}{km} suhunya
$20 + 75 = \qty{95}{\degreeCelsius}$ sampai $20 + 102 = \qty{122}{\degreeCelsius}$:
di dalam \omterm{def:g11:fossil-fuel-geology:oil-window}{jendela minyak}, atau di dasarnya untuk gradien yang paling curam.
Aturan itu masuk akal: di kebanyakan cekungan, \omterm{def:g11:fossil-fuel-geology:source-rock}{batuan sumber} di bawah
\qty{3}{km} batuan telah mulai menghasilkan minyak bumi.
\end{solution}

\begin{solution}{exo:g11:fossil-fuel-geology:11}
(1) \qty{30}{\percent}: lignit; \qty{70}{\percent}: \omterm{def:g11:fossil-fuel-geology:coal}{batu bara} bituminus;
\qty{92}{\percent}: antrasit. (2) Antrasit, di luar sekitar
\qty{5}{km} dan \qty{150}{\degreeCelsius}. (3) Hitam dan berkilap kaca.
\end{solution}

\begin{solution}{exo:g11:fossil-fuel-geology:12}
(1) $1100/25 = \qty{44}{m}$ per sejuta tahun. (2) $0.044 \times 30
\approx \qty{1.3}{\degreeCelsius}$ per sejuta tahun. (3) Jendela itu
membentang \qty{60}{\degreeCelsius}: $60/1.32 \approx 45$ juta tahun.
(4) Penimbunan melambat setelah 95 juta tahun yang lalu (sekitar 24, lalu 23
dan \qty{20}{m} per sejuta tahun): \omterm{def:g11:fossil-fuel-geology:source-rock}{batuan sumber} menghangat lebih lambat.
\end{solution}

\begin{solution}{exo:g11:fossil-fuel-geology:13}
(1) $(80 - 51)/(117 - 51) = 29/66 \approx 0.44$: sekitar \qty{44}{\percent}.
(2) Minyak bumi itu bermigrasi ke atas sebelum ada \omterm{def:g11:fossil-fuel-geology:trap}{perangkap} di atasnya: minyak
itu lolos ke permukaan, tempat ia teroksidasi, atau berkumpul di \omterm{def:g11:fossil-fuel-geology:trap}{perangkap}
lain. (3) Pada 40 juta tahun yang lalu \omterm{def:g11:fossil-fuel-geology:source-rock}{batuan sumber} telah meninggalkan \omterm{def:g11:fossil-fuel-geology:oil-window}{jendela
minyak}: batuan itu hanya menghasilkan gas, dan antiklin itu hanya dapat menahan
gas darinya.
\end{solution}

\begin{solution}{exo:g11:fossil-fuel-geology:14}
(1) Puncak: $(60 - 20)/40 = \qty{1.0}{km}$; dasar: $(120 - 20)/40 =
\qty{2.5}{km}$. (2) \qty{1.0}{km} terletak antara \qty{0.4}{km} (150 juta
tahun yang lalu) dan \qty{1.2}{km} (120): $150 - 30 \times 0.6/0.8 = 127.5$
juta tahun yang lalu. \qty{2.5}{km} terletak antara \qty{2.3}{km} (95) dan
\qty{3.0}{km} (66): $95 - 29 \times 0.2/0.7 \approx 86.7$ juta tahun
yang lalu. (3) Tidak ada minyak bumi: \omterm{def:g11:fossil-fuel-geology:source-rock}{batuan sumber} meninggalkan \omterm{def:g11:fossil-fuel-geology:oil-window}{jendela minyak}
sebelum antiklin itu terbentuk, dan minyak buminya telah bermigrasi pergi.
Antiklin itu dapat menahan gas yang terbentuk setelah 80 juta tahun yang lalu.
\end{solution}

\begin{solution}{exo:g11:fossil-fuel-geology:15}
(1) Setidaknya \qty{5}{km}. (2) $15 + 30 \times 5 = \qty{165}{\degreeCelsius}$,
di atas \qty{150}{\degreeCelsius}: sesuai. (3) Batuan setebal setidaknya
\qty{5}{km} telah terangkat dan tererosi. $\qty{5e6}{mm}$ dalam paling lama
\num{299} juta tahun memberikan setidaknya $5 \times 10^6/2.99 \times 10^8
\approx \qty{0.017}{mm}$ setahun secara rata-rata.
\end{solution}

\begin{solution}{pb:g11:fossil-fuel-geology:1}
\textbf{1.} \omterm{def:g11:fossil-fuel-geology:source-rock}{Batuan sumber}: serpih hitam; reservoir: batu pasir berpori; \omterm{def:g11:fossil-fuel-geology:reservoir}{batuan
penutup}: batu lumpur di atasnya.

\textbf{2.} Plankton yang mengendap di dasar laut yang miskin oksigen luput dari
pembusukan dan terkubur di dalam lumpur.

\textbf{3.} \omterm{def:g11:fossil-fuel-geology:source-rock}{Batuan sumber} pada zaman Jura (165 terletak antara 201.4 dan
143.1). Batu pasirnya mulai diendapkan pada zaman Jura dan sebagian besar
diendapkan pada Kapur Awal.

\textbf{4.} Bahan organik yang tersebar dalam batuan sedimen, yang tidak larut
dalam pelarut organik. \omterm{def:g11:fossil-fuel-geology:source-rock}{Kerogen} baru menghasilkan minyak bumi setelah
terpanaskan sampai sekitar \qty{60}{\degreeCelsius}, dan itu memerlukan
penimbunan.

\textbf{5.} $400 + 800 + 1100 + 700 + 700 + 300 - 500 = \qty{3500}{m}$.

\textbf{6.} \qty{0.4}{km}, \qty{1.2}{km}, \qty{2.3}{km}, \qty{3.0}{km},
\qty{3.7}{km}, \qty{4.0}{km}; kini \qty{3.5}{km}.

\textbf{7.} $400/15 \approx 26.7$; $800/30 \approx 26.7$; $1100/25 = 44$;
$700/29 \approx 24.1$; $700/31 \approx 22.6$; $300/15 = 20$ (\unit{m} per
sejuta tahun).

\textbf{8.} 20 juta tahun yang lalu, pada \qty{4.0}{km}.

\textbf{9.} Cekungan itu terangkat dan batuan setebal \qty{500}{m} tererosi
dari bagian atasnya.

\textbf{10.} Dasar batu pasir diendapkan 150 juta tahun yang lalu, di atas
serpih dan batu lumpur setebal \qty{400}{m}: $3.5 - 0.4 = \qty{3.1}{km}$.

\textbf{11.} $T = 20 + 30 z$: \qty{32}{\degreeCelsius}, \num{56}, \num{89},
\num{110}, \num{131}, \qty{140}{\degreeCelsius}; kini
\qty{125}{\degreeCelsius}.

\textbf{12.} $40/30 \approx \qty{1.33}{km}$ dan $100/30 \approx
\qty{3.33}{km}$.

\textbf{13.} Antara 120 juta tahun yang lalu (\qty{56}{\degreeCelsius}) dan
95 juta tahun yang lalu (\qty{89}{\degreeCelsius}).

\textbf{14.} Suhunya naik 33 derajat dalam 25 juta tahun; diperlukan
$25 \times 4/33 \approx 3.0$ juta tahun untuk naik 4 derajat dari
\num{56} ke \qty{60}{\degreeCelsius}: sekitar $120 - 3.0 = 117$ juta tahun
yang lalu.

\textbf{15.} Antara 66 juta tahun yang lalu (\qty{110}{\degreeCelsius}) dan
35 (\qty{131}{\degreeCelsius}): $66 - 31 \times 10/21 \approx 51.2$ juta
tahun yang lalu.

\textbf{16.} \qty{140}{\degreeCelsius}, 20 juta tahun yang lalu. Pada
\qty{125}{\degreeCelsius} kini batuan itu berada di dalam \omterm{def:g11:fossil-fuel-geology:oil-window}{jendela gas}: batuan
itu terutama menghasilkan metana.

\textbf{17.} Karena tidak ada \omterm{def:g11:fossil-fuel-geology:trap}{perangkap} di atasnya, minyak bumi itu bermigrasi
ke atas dan lolos, ke permukaan atau ke \omterm{def:g11:fossil-fuel-geology:trap}{perangkap} lain.

\textbf{18.} $(80 - 51.2)/(117.0 - 51.2) \approx 0.44$.

\textbf{19.} Gas di puncak, lalu minyak bumi, lalu air.

\textbf{20.} \omterm{def:g11:fossil-fuel-geology:source-rock}{Batuan sumber} memasuki \omterm{def:g11:fossil-fuel-geology:oil-window}{jendela minyak} sekitar \textbf{117
juta tahun yang lalu}.
\end{solution}
